The simplest version of the transition matrix considers only how a single roll of two six-sided dice moves the player around the 40 spaces in $S$. The sum of two six-sided dice for each roll can be between 2 and 12. Table~\ref{tab:dice} shows the probability of each sum.

\begin{table}[H]
\centering
\begin{tabular}{|c|l|c|}
\hline
\textbf{Sum} & \textbf{Dice Roll Combinations} & \textbf{Probability} \\
\hline
2  & (1,1) & 1/36 \\
\hline
3  & (1,2), (2,1) & 2/36 \\
\hline
4  & (1,3), (2,2), (3,1) & 3/36 \\
\hline
5  & (1,4), (2,3), (3,2), (4,1) & 4/36 \\
\hline
6  & (1,5), (2,4), (3,3), (4,2), (5,1) & 5/36 \\
\hline
7  & (1,6), (2,5), (3,4), (4,3), (5,2), (6,1) & 6/36 \\
\hline
8  & (2,6), (3,5), (4,4), (5,3), (6,2) & 5/36 \\
\hline
9  & (3,6), (4,5), (5,4), (6,3) & 4/36 \\
\hline
10 & (4,6), (5,5), (6,4) & 3/36 \\
\hline
11 & (5,6), (6,5) & 2/36 \\
\hline
12 & (6,6) & 1/36 \\
\hline
\end{tabular}
\caption{Dice roll combinations and probability of each sum for two six-sided dice}
\label{tab:dice}
\end{table}

